\documentclass[10pt,a4paper]{article}
\usepackage[margin=0.8cm]{geometry}\usepackage{fontspec}\usepackage{tabularx}
\usepackage{array}\usepackage{booktabs}\usepackage{fancyhdr}\usepackage{setspace}
\usepackage{amssymb}\usepackage{xcolor}
\setmainfont{Times New Roman}
\newfontfamily\hindifont{Nirmala UI}
\pagestyle{fancy}\fancyhf{}
\fancyhead[L]{\small\textbf{AIGGPA Research Schedule}}
\fancyhead[R]{\small Respondent ID: \_\_\_\_\_\_\_\_}
\fancyfoot[C]{\small Page \thepage}\renewcommand{\headrulewidth}{0.5pt}
\setlength{\headheight}{14pt}\setstretch{1.05}
\newcolumntype{L}[1]{>{\raggedright\arraybackslash}p{#1}}
\newcommand{\likert}{{\Large$\square$}\,1\quad{\Large$\square$}\,2\quad{\Large$\square$}\,3\quad{\Large$\square$}\,4\quad{\Large$\square$}\,5}
\newcommand{\ymark}{{\Large$\square$}\,Yes/{\hindifont\itshape हाँ}\quad{\Large$\square$}\,No/{\hindifont\itshape नहीं}}
\newcommand{\blank}{\rule{7cm}{0.4pt}}
\begin{document}

\begin{center}
{\Large\bfseries AIGGPA FIELD RESEARCH SCHEDULE}\\[1pt]
{\large\hindifont\itshape\color{black!70} AIGGPA \textnormal{क्षेत्र अनुसंधान अनुसूची}}\\[6pt]
{\small Assessment of Digital Tool Usage Among Government Employees}\\[1pt]
{\small\hindifont\itshape\color{black!70} सरकारी कर्मचारियों में डिजिटल उपकरण उपयोग का मूल्यांकन}\\[4pt]
{\small AIGGPA Bhopal \quad$\bullet$\quad 2026}
\end{center}
\vspace{4pt}\noindent\rule{\textwidth}{1pt}\vspace{4pt}

\noindent\begin{tabularx}{\textwidth}{@{}L{9.5cm} X@{}}
\textbf{Respondent ID:} \blank & \textbf{Date:} \_\_\_/\_\_\_/2026 \\[6pt]
\textbf{Department:} {\Large$\square$} Revenue {\Large$\square$} Rural Dev {\Large$\square$} Forest {\Large$\square$} Health & \textbf{Office:} {\Large$\square$} HO {\Large$\square$} District \\[6pt]
\textbf{Interviewer:} \blank & \textbf{Time:} \_\_:\_\_\quad to \quad\_\_:\_\_ \\
\end{tabularx}
\vspace{4pt}\noindent\rule{\textwidth}{0.5pt}\vspace{4pt}
\noindent\rule{\textwidth}{0.3pt}\vspace{3pt}
\noindent{\large\bfseries Section A: Respondent Profile}\\
{\hindifont\itshape\color{black!70} उत्तरदाता प्रोफाइल}\vspace{3pt}

\noindent\begin{tabularx}{\textwidth}{@{}L{0.8cm} L{7.7cm} X@{}}
Q1 & Designation / Post: / {\hindifont\small\itshape\color{black!70} पदनाम / पद:} & \blank \\[3pt]
Q2 & Job Role / Level: / {\hindifont\small\itshape\color{black!70} कार्य भूमिका / स्तर:} & \blank \\[3pt]
Q3 & Age group: / {\hindifont\small\itshape\color{black!70} आयु वर्ग:} & {\Large$\square$} Below 30/{\hindifont 30 से कम}\quad{\Large$\square$} 30--45\quad{\Large$\square$} 46--60 \\[3pt]
Q4 & Gender: / {\hindifont\small\itshape\color{black!70} लिंग:} & {\Large$\square$} Male/{\hindifont पुरुष}\quad{\Large$\square$} Female/{\hindifont महिला}\quad{\Large$\square$} Other/{\hindifont अन्य} \\[3pt]
Q5 & Years of service: / {\hindifont\small\itshape\color{black!70} सेवा के वर्ष:} & {\Large$\square$} 0--5\quad{\Large$\square$} 6--10\quad{\Large$\square$} 11--20\quad{\Large$\square$} 21+ \\[3pt]
Q6 & Highest education: / {\hindifont\small\itshape\color{black!70} उच्चतम शिक्षा:} & {\Large$\square$} Up to 12th/{\hindifont\small\itshape\color{black!70} 12वीं तक}\quad{\Large$\square$} Grad/{\hindifont\small\itshape\color{black!70} स्नातक}\quad{\Large$\square$} PG/{\hindifont\small\itshape\color{black!70} स्नातकोत्तर}\quad{\Large$\square$} Prof/{\hindifont\small\itshape\color{black!70} व्यावसायिक} \\[3pt]
\end{tabularx}
\vspace{4pt}
\noindent\rule{\textwidth}{0.3pt}\vspace{3pt}
\noindent{\large\bfseries Section B: Infrastructure & Facilitating Conditions}\\
{\hindifont\itshape\color{black!70} अवसंरचना एवं सुविधाएँ}\vspace{3pt}

\small\textit{Likert: 1=Strongly Disagree \quad 2=Disagree \quad 3=Neutral \quad 4=Agree \quad 5=Strongly Agree}
\\
{\hindifont\small\textit{लिकर्ट: 1=पूर्णतः असहमत \quad 2=असहमत \quad 3=तटस्थ \quad 4=सहमत \quad 5=पूर्णतः सहमत}}\vspace{3pt}

\noindent\begin{tabularx}{\textwidth}{@{}L{0.8cm} X L{5.2cm}@{}}
Q7 & What digital devices are available at your workstation? & {\Large$\square$} Desktop {\Large$\square$} Laptop {\Large$\square$} Tablet {\Large$\square$} Phone {\Large$\square$} None/{\hindifont कोई नहीं} \\
 & {\hindifont\small\itshape\color{black!70} आपके कार्यस्थल पर कौन-से डिजिटल उपकरण उपलब्ध हैं?} & \\[3pt]
Q8 & Do you share your device with other employees? & {\Large$\square$} Yes, always/{\hindifont हमेशा} {\Large$\square$} Sometimes/{\hindifont कभी-कभी} {\Large$\square$} No, dedicated/{\hindifont नहीं} \\
 & {\hindifont\small\itshape\color{black!70} क्या आप अपना उपकरण अन्य कर्मचारियों के साथ साझा करते हैं?} & \\[3pt]
Q9 & Rate internet connectivity at your office: & \likert \\
 & {\hindifont\small\itshape\color{black!70} अपने कार्यालय में इंटरनेट कनेक्टिविटी का मूल्यांकन करें:} & \\[3pt]
Q10 & How often do you experience internet outages per week? & {\Large$\square$} Never {\Large$\square$} 1--2 {\Large$\square$} 3--5 {\Large$\square$} Daily/{\hindifont प्रतिदिन} \\
 & {\hindifont\small\itshape\color{black!70} सप्ताह में कितनी बार इंटरनेट बंद होता है?} & \\[3pt]
Q11 & Is IT helpdesk / technical support available? & \ymark \\
 & {\hindifont\small\itshape\color{black!70} क्या IT हेल्पडेस्क / तकनीकी सहायता उपलब्ध है?} & \\[3pt]
Q12 & If yes, how quickly are issues resolved? & {\Large$\square$} Same day/{\hindifont उसी दिन} {\Large$\square$} 2--3 days {\Large$\square$} 1 week+ {\Large$\square$} Never/{\hindifont कभी नहीं} \\
 & {\hindifont\small\itshape\color{black!70} यदि हाँ, तो समस्याएँ कितनी जल्दी हल होती हैं?} & \\[3pt]
\end{tabularx}
\vspace{4pt}
\noindent\rule{\textwidth}{0.3pt}\vspace{3pt}
\noindent{\large\bfseries Section C: Performance Expectancy}\\
{\hindifont\itshape\color{black!70} प्रदर्शन अपेक्षा}\vspace{3pt}

\noindent\begin{tabularx}{\textwidth}{@{}L{0.8cm} X L{5.2cm}@{}}
Q13 & Digital tools help me complete tasks faster than paper. & \likert \\
 & {\hindifont\small\itshape\color{black!70} डिजिटल उपकरण मुझे कागज़ की तुलना में कार्य तेज़ी से पूरा करने में मदद करते हैं।} & \\[3pt]
Q14 & Digital tools improve the quality/accuracy of my work. & \likert \\
 & {\hindifont\small\itshape\color{black!70} डिजिटल उपकरण मेरे काम की गुणवत्ता/सटीकता में सुधार करते हैं।} & \\[3pt]
Q15 & Using digital tools increases my overall productivity. & \likert \\
 & {\hindifont\small\itshape\color{black!70} डिजिटल उपकरणों का उपयोग मेरी समग्र उत्पादकता बढ़ाता है।} & \\[3pt]
Q16 & The digital tools available are well-suited to my actual job tasks. & \likert \\
 & {\hindifont\small\itshape\color{black!70} उपलब्ध डिजिटल उपकरण मेरे वास्तविक कार्यों के लिए उपयुक्त हैं।} & \\[3pt]
\end{tabularx}
\vspace{4pt}
\noindent\rule{\textwidth}{0.3pt}\vspace{3pt}
\noindent{\large\bfseries Section D: Effort Expectancy}\\
{\hindifont\itshape\color{black!70} प्रयास अपेक्षा}\vspace{3pt}

\noindent\begin{tabularx}{\textwidth}{@{}L{0.8cm} X L{5.2cm}@{}}
Q17 & How difficult do you find digital tools to use? & \likert{} {\scriptsize(1=Very Easy, 5=Very Difficult)} \\
 & {\hindifont\small\itshape\color{black!70} डिजिटल उपकरणों का उपयोग आपको कितना कठिन लगता है?} & \\[3pt]
Q18 & I am confident in my ability to use digital tools for my work. & \likert \\
 & {\hindifont\small\itshape\color{black!70} मुझे डिजिटल उपकरण उपयोग करने की अपनी क्षमता पर विश्वास है।} & \\[3pt]
Q19 & Learning to use a new portal/app takes me: & {\Large$\square$} <1 day {\Large$\square$} Few days {\Large$\square$} 1--2 weeks {\Large$\square$} >2 weeks \\
 & {\hindifont\small\itshape\color{black!70} एक नया पोर्टल/ऐप सीखने में मुझे लगता है:} & \\[3pt]
Q20 & The design/interface of most government portals is user-friendly. & \likert \\
 & {\hindifont\small\itshape\color{black!70} अधिकांश सरकारी पोर्टलों का डिज़ाइन उपयोगकर्ता-अनुकूल है।} & \\[3pt]
\end{tabularx}
\vspace{4pt}
\noindent\rule{\textwidth}{0.3pt}\vspace{3pt}
\noindent{\large\bfseries Section E: Social Influence}\\
{\hindifont\itshape\color{black!70} सामाजिक प्रभाव}\vspace{3pt}

\noindent\begin{tabularx}{\textwidth}{@{}L{0.8cm} X L{5.2cm}@{}}
Q21 & My superiors encourage the use of digital tools. & \likert \\
 & {\hindifont\small\itshape\color{black!70} मेरे वरिष्ठ अधिकारी डिजिटल उपकरणों के उपयोग को प्रोत्साहित करते हैं।} & \\[3pt]
Q22 & My colleagues regularly use digital tools in their work. & \likert \\
 & {\hindifont\small\itshape\color{black!70} मेरे सहकर्मी नियमित रूप से डिजिटल उपकरणों का उपयोग करते हैं।} & \\[3pt]
Q23 & There is a formal mandate/order requiring digital tool use in my department. & {\Large$\square$} Yes/{\hindifont हाँ} {\Large$\square$} No/{\hindifont नहीं} {\Large$\square$} Don't know/{\hindifont पता नहीं} \\
 & {\hindifont\small\itshape\color{black!70} मेरे विभाग में डिजिटल उपकरण के लिए औपचारिक आदेश है।} & \\[3pt]
\end{tabularx}
\vspace{4pt}
\noindent\rule{\textwidth}{0.3pt}\vspace{3pt}
\noindent{\large\bfseries Section F: Awareness \& Usage}\\
{\hindifont\itshape\color{black!70} जागरूकता एवं उपयोग}\vspace{3pt}

\noindent\begin{tabularx}{\textwidth}{@{}L{0.8cm} X L{5.2cm}@{}}
Q24 & Which general tools are you aware of? & {\Large$\square$} e-Office {\Large$\square$} CM Helpline {\Large$\square$} PFMS {\Large$\square$} SPARROW {\Large$\square$} iGOT {\Large$\square$} MP eDistrict \\
 & {\hindifont\small\itshape\color{black!70} आप किन सामान्य उपकरणों से अवगत हैं?} & \\[3pt]
Q25 & How often do you use digital tools for work? & {\Large$\square$} Daily/{\hindifont प्रतिदिन} {\Large$\square$} Weekly/{\hindifont साप्ताहिक} {\Large$\square$} Monthly/{\hindifont मासिक} {\Large$\square$} Rarely/{\hindifont कभी-कभी} {\Large$\square$} Never/{\hindifont कभी नहीं} \\
 & {\hindifont\small\itshape\color{black!70} आप कितनी बार काम के लिए डिजिटल उपकरणों का उपयोग करते हैं?} & \\[3pt]
Q26 & What percentage of your work is done digitally? & {\Large$\square$} 0--20\% {\Large$\square$} 21--40\% {\Large$\square$} 41--60\% {\Large$\square$} 61--80\% {\Large$\square$} 81--100\% \\
 & {\hindifont\small\itshape\color{black!70} आपके कार्य का कितना प्रतिशत डिजिटल रूप से होता है?} & \\[3pt]
\end{tabularx}
\newpage
\noindent\rule{\textwidth}{0.3pt}\vspace{3pt}
\noindent{\large\bfseries Section G: Training \& Capacity Building}\\
{\hindifont\itshape\color{black!70} प्रशिक्षण एवं क्षमता निर्माण}\vspace{3pt}

\noindent\begin{tabularx}{\textwidth}{@{}L{0.8cm} X L{5.2cm}@{}}
Q27 & Have you attended any digital skills training in the last 2 years? & \ymark \\
 & {\hindifont\small\itshape\color{black!70} क्या पिछले 2 वर्षों में आपने कोई डिजिटल कौशल प्रशिक्षण लिया है?} & \\[3pt]
Q28 & If yes, how many training sessions? & {\Large$\square$} 1 {\Large$\square$} 2--3 {\Large$\square$} 4--5 {\Large$\square$} More than 5 \\
 & {\hindifont\small\itshape\color{black!70} यदि हाँ, तो कितने प्रशिक्षण सत्र?} & \\[3pt]
Q29 & Rate the quality of training received: & \likert \\
 & {\hindifont\small\itshape\color{black!70} प्राप्त प्रशिक्षण की गुणवत्ता का मूल्यांकन करें:} & \\[3pt]
Q30 & Was the training sufficient for your actual job needs? & \likert \\
 & {\hindifont\small\itshape\color{black!70} क्या प्रशिक्षण आपकी वास्तविक कार्य आवश्यकताओं के लिए पर्याप्त था?} & \\[3pt]
Q31 & What topics need more training? & \rule{\linewidth}{0.4pt}\\[2pt] & & \rule{\linewidth}{0.4pt} \\
 & {\hindifont\small\itshape\color{black!70} किन विषयों में अधिक प्रशिक्षण चाहिए?} & \\[3pt]
\end{tabularx}
\vspace{4pt}
\noindent\rule{\textwidth}{0.3pt}\vspace{3pt}
\noindent{\large\bfseries Section H: Challenges \& Barriers}\\
{\hindifont\itshape\color{black!70} चुनौतियाँ एवं बाधाएँ}\vspace{3pt}

\noindent\begin{tabularx}{\textwidth}{@{}L{0.8cm} X L{5.2cm}@{}}
Q32 & What issues do you face? (tick all) & {\Large$\square$} Slow internet {\Large$\square$} Crashes {\Large$\square$} No device {\Large$\square$} Complex UI {\Large$\square$} No training {\Large$\square$} No support {\Large$\square$} Power cuts \\
 & {\hindifont\small\itshape\color{black!70} आपको क्या समस्याएँ आती हैं? (सभी पर टिक करें)} & \\[3pt]
Q33 & How often do digital issues disrupt your work? & {\Large$\square$} Daily {\Large$\square$} Weekly {\Large$\square$} Monthly {\Large$\square$} Rarely {\Large$\square$} Never \\
 & {\hindifont\small\itshape\color{black!70} डिजिटल समस्याएँ कितनी बार आपके काम में बाधा डालती हैं?} & \\[3pt]
Q34 & I feel comfortable asking for help with digital tools. & \likert \\
 & {\hindifont\small\itshape\color{black!70} मुझे डिजिटल उपकरणों के लिए सहायता माँगने में सहजता है।} & \\[3pt]
Q35 & My organisation provides adequate support for digital tools. & \likert \\
 & {\hindifont\small\itshape\color{black!70} मेरा संगठन डिजिटल उपकरणों के लिए पर्याप्त सहायता प्रदान करता है।} & \\[3pt]
Q36 & My department is committed to digital transformation. & \likert \\
 & {\hindifont\small\itshape\color{black!70} मेरा विभाग डिजिटल परिवर्तन के प्रति प्रतिबद्ध है।} & \\[3pt]
\end{tabularx}
\vspace{4pt}
\noindent\rule{\textwidth}{0.3pt}\vspace{3pt}
\noindent{\large\bfseries Section I: Recommendations}\\
{\hindifont\itshape\color{black!70} सुझाव}\vspace{3pt}

\noindent Q37. Rank these priorities (1=highest, 5=lowest): / {\hindifont\small\itshape\color{black!70} इन प्राथमिकताओं को क्रम दें (1=सर्वोच्च, 5=न्यूनतम):}\vspace{3pt}

\noindent\begin{tabularx}{\textwidth}{@{}L{0.8cm} X L{1.5cm}@{}}
 & Better internet / {\hindifont\small\itshape\color{black!70} बेहतर इंटरनेट} & Rank: \_\_\_ \\[3pt]
 & More/better devices / {\hindifont\small\itshape\color{black!70} अधिक/बेहतर उपकरण} & Rank: \_\_\_ \\[3pt]
 & More training / {\hindifont\small\itshape\color{black!70} अधिक प्रशिक्षण} & Rank: \_\_\_ \\[3pt]
 & Simpler portals / {\hindifont\small\itshape\color{black!70} सरल पोर्टल} & Rank: \_\_\_ \\[3pt]
 & Faster IT support / {\hindifont\small\itshape\color{black!70} तेज़ IT सहायता} & Rank: \_\_\_ \\[3pt]
\end{tabularx}
\vspace{4pt}
\noindent Q38. What one change would most improve your use of digital tools?\\
{\hindifont\small\itshape\color{black!70} कौन-सा एक बदलाव आपके डिजिटल उपकरण उपयोग में सबसे अधिक सुधार करेगा?}\vspace{3pt}

\noindent\rule{\textwidth}{0.4pt}\vspace{4pt}
\noindent\rule{\textwidth}{0.4pt}\vspace{4pt}
\noindent Q39. Has digital tool adoption improved service delivery to citizens?\\
{\hindifont\small\itshape\color{black!70} क्या डिजिटल उपकरणों से नागरिकों को सेवा वितरण में सुधार हुआ है?}\vspace{2pt}

\noindent {\Large$\square$} Yes, significantly/{\hindifont हाँ, काफ़ी}\quad{\Large$\square$} Somewhat/{\hindifont कुछ हद तक}\quad{\Large$\square$} No change/{\hindifont कोई बदलाव नहीं}\quad{\Large$\square$} Worsened/{\hindifont बिगड़ा}\quad{\Large$\square$} Can't say/{\hindifont कह नहीं सकते}

\vspace{4pt}\noindent\rule{\textwidth}{1pt}\vspace{2pt}
\begin{center}\textit{--- End of Common Schedule (Q1--Q39) / {\hindifont सामान्य अनुसूची समाप्त} ---}\end{center}
\rule{\textwidth}{0.5pt}
\newpage
\noindent\rule{\textwidth}{0.3pt}\vspace{3pt}
\noindent{\large\bfseries Section J: Department-Specific --- REVENUE / {\hindifont राजस्व विभाग}}\vspace{3pt}

\small\textit{Administer ONLY to Revenue respondents / {\hindifont केवल राजस्व विभाग के उत्तरदाताओं के लिए}}\vspace{3pt}

\noindent\begin{tabularx}{\textwidth}{@{}L{0.8cm} X L{5.2cm}@{}}
Q40 & Which Revenue digital tools are you aware of? & {\Large$\square$} Bhulekh/WebGIS {\Large$\square$} RCMS {\Large$\square$} SAARA {\Large$\square$} SAMPADA {\Large$\square$} e-Court {\Large$\square$} None/{\hindifont कोई नहीं} \\
 & {\hindifont\small\itshape\color{black!70} राजस्व विभाग के कौन-से डिजिटल उपकरणों से आप अवगत हैं?} & \\[3pt]
Q41 & Has Bhulekh/WebGIS improved speed and accuracy of land record verification? & \likert \\
 & {\hindifont\small\itshape\color{black!70} क्या भूलेख/WebGIS ने भूमि अभिलेख सत्यापन की गति और सटीकता में सुधार किया है?} & \\[3pt]
Q42 & How difficult is RCMS to use for tracking revenue cases? & \likert{} {\scriptsize(1=Easy, 5=Difficult)} \\
 & {\hindifont\small\itshape\color{black!70} राजस्व मामलों की ट्रैकिंग के लिए RCMS का उपयोग कितना कठिन है?} & \\[3pt]
Q43 & What percentage of land records still need physical paper files? & {\Large$\square$} 0--20\% {\Large$\square$} 21--40\% {\Large$\square$} 41--60\% {\Large$\square$} 61--80\% {\Large$\square$} 81--100\% \\
 & {\hindifont\small\itshape\color{black!70} भूमि अभिलेखों का कितना प्रतिशत अभी भी कागज़ी फाइलों की ज़रूरत है?} & \\[3pt]
Q44 & How often do citizens visit expecting services through digital tools? & \likert{} {\scriptsize(1=Never, 5=Daily)} \\
 & {\hindifont\small\itshape\color{black!70} नागरिक कितनी बार डिजिटल सेवाओं की अपेक्षा से आते हैं?} & \\[3pt]
\end{tabularx}
\vspace{3pt}
\noindent\rule{\textwidth}{0.3pt}\vspace{3pt}
\noindent{\large\bfseries Section K: Department-Specific --- RURAL DEVELOPMENT / {\hindifont ग्रामीण विकास}}\vspace{3pt}

\small\textit{Administer ONLY to Rural Development respondents / {\hindifont केवल ग्रामीण विकास उत्तरदाताओं के लिए}}\vspace{3pt}

\noindent\begin{tabularx}{\textwidth}{@{}L{0.8cm} X L{5.2cm}@{}}
Q45 & Which Rural Development tools are you aware of? & {\Large$\square$} NREGASoft/NMMS {\Large$\square$} e-Gram Swaraj {\Large$\square$} PMAY-G {\Large$\square$} SBM-G {\Large$\square$} Panchayat Darpan {\Large$\square$} PFMS {\Large$\square$} None \\
 & {\hindifont\small\itshape\color{black!70} ग्रामीण विकास के कौन-से उपकरणों से आप अवगत हैं?} & \\[3pt]
Q46 & How difficult is managing multiple portals at the same time? & \likert{} {\scriptsize(1=Easy, 5=Difficult)} \\
 & {\hindifont\small\itshape\color{black!70} एक साथ कई पोर्टलों का प्रबंधन कितना कठिन है?} & \\[3pt]
Q47 & Rate internet connectivity at block/panchayat offices: & \likert \\
 & {\hindifont\small\itshape\color{black!70} ब्लॉक/पंचायत कार्यालयों में इंटरनेट कनेक्टिविटी का मूल्यांकन करें:} & \\[3pt]
Q48 & Has the NMMS app improved accuracy of MGNREGA attendance? & \likert \\
 & {\hindifont\small\itshape\color{black!70} क्या NMMS ऐप ने मनरेगा उपस्थिति की सटीकता में सुधार किया है?} & \\[3pt]
Q49 & How much of your working day is spent entering data into portals? & {\Large$\square$} <1hr {\Large$\square$} 1--2hr {\Large$\square$} 2--4hr {\Large$\square$} 4+hr {\Large$\square$} Almost all day \\
 & {\hindifont\small\itshape\color{black!70} आपके कार्यदिवस का कितना समय पोर्टलों में डेटा दर्ज करने में लगता है?} & \\[3pt]
\end{tabularx}
\vspace{3pt}
\noindent\rule{\textwidth}{0.3pt}\vspace{3pt}
\noindent{\large\bfseries Section L: Department-Specific --- FOREST / {\hindifont वन विभाग}}\vspace{3pt}

\small\textit{Administer ONLY to Forest respondents / {\hindifont केवल वन विभाग उत्तरदाताओं के लिए}}\vspace{3pt}

\noindent\begin{tabularx}{\textwidth}{@{}L{0.8cm} X L{5.2cm}@{}}
Q50 & Which Forest department digital tools are you aware of? & {\Large$\square$} e-Green Watch {\Large$\square$} AI Alert {\Large$\square$} GIS {\Large$\square$} Forest Offence MIS {\Large$\square$} Nursery MIS {\Large$\square$} None \\
 & {\hindifont\small\itshape\color{black!70} वन विभाग के कौन-से डिजिटल उपकरणों से आप अवगत हैं?} & \\[3pt]
Q51 & Has the AI alert system improved detection of illegal activity? & \likert{}\quad{\Large$\square$} N/A \\
 & {\hindifont\small\itshape\color{black!70} क्या AI अलर्ट प्रणाली ने अवैध गतिविधि की पहचान में सुधार किया है?} & \\[3pt]
Q52 & How difficult do you find GIS tools? & \likert{} {\scriptsize(1=Easy, 5=Difficult)} \\
 & {\hindifont\small\itshape\color{black!70} GIS उपकरण आपको कितने कठिन लगते हैं?} & \\[3pt]
Q53 & Do you have a GPS-enabled device for field verification? & {\Large$\square$} Dept-issued {\Large$\square$} Personal device {\Large$\square$} Not available/{\hindifont उपलब्ध नहीं} \\
 & {\hindifont\small\itshape\color{black!70} क्या आपके पास क्षेत्र सत्यापन के लिए GPS उपकरण है?} & \\[3pt]
Q54 & How often does poor network prevent digital tool use in the field? & \likert{} {\scriptsize(1=Never, 5=Always)} \\
 & {\hindifont\small\itshape\color{black!70} खराब नेटवर्क कितनी बार क्षेत्र में डिजिटल उपकरण उपयोग रोकता है?} & \\[3pt]
\end{tabularx}
\vspace{3pt}
\noindent\rule{\textwidth}{0.3pt}\vspace{3pt}
\noindent{\large\bfseries Section M: Department-Specific --- HEALTH / {\hindifont स्वास्थ्य विभाग}}\vspace{3pt}

\small\textit{Administer ONLY to Health respondents / {\hindifont केवल स्वास्थ्य विभाग उत्तरदाताओं के लिए}}\vspace{3pt}

\noindent\begin{tabularx}{\textwidth}{@{}L{0.8cm} X L{5.2cm}@{}}
Q55 & Which Health department digital tools are you aware of? & {\Large$\square$} ANMOL {\Large$\square$} HMIS {\Large$\square$} Nikshay {\Large$\square$} eVIN {\Large$\square$} IHIP {\Large$\square$} ABHA {\Large$\square$} MPCDSR {\Large$\square$} None \\
 & {\hindifont\small\itshape\color{black!70} स्वास्थ्य विभाग के कौन-से डिजिटल उपकरणों से आप अवगत हैं?} & \\[3pt]
Q56 & Has ANMOL MP or ABHA improved your ability to track patients? & \likert \\
 & {\hindifont\small\itshape\color{black!70} क्या ANMOL या ABHA ने मरीज़ ट्रैकिंग में सुधार किया है?} & \\[3pt]
Q57 & How often do you enter the SAME data in both paper AND digital? & \likert{} {\scriptsize(1=Never, 5=Always)} \\
 & {\hindifont\small\itshape\color{black!70} कितनी बार आप एक ही डेटा कागज़ और डिजिटल दोनों में दर्ज करते हैं?} & \\[3pt]
Q58 & How reliably does eVIN reflect actual vaccine stock? & \likert{} {\scriptsize(1=Unreliable, 5=Reliable)} \\
 & {\hindifont\small\itshape\color{black!70} eVIN वास्तविक टीका स्टॉक को कितनी विश्वसनीयता से दर्शाता है?} & \\[3pt]
Q59 & How much does mandatory IHIP disease reporting add to your workload? & \likert{} {\scriptsize(1=None, 5=Significant)} \\
 & {\hindifont\small\itshape\color{black!70} अनिवार्य IHIP रोग रिपोर्टिंग आपके कार्यभार में कितना जोड़ती है?} & \\[3pt]
\end{tabularx}
\newpage
\noindent\rule{\textwidth}{0.3pt}\vspace{3pt}
\noindent{\large\bfseries Section N: Role-Specific Questions}\\
{\hindifont\itshape\color{black!70} भूमिका-विशिष्ट प्रश्न}\vspace{3pt}

\noindent\begin{tabularx}{\textwidth}{@{}L{0.8cm} X L{5.2cm}@{}}
Q60 & Primary interaction with digital tools: & {\Large$\square$} Data entry {\Large$\square$} Review/approve {\Large$\square$} Field verification {\Large$\square$} Don't use {\Large$\square$} Other \\
 & {\hindifont\small\itshape\color{black!70} डिजिटल उपकरणों के साथ आपकी प्राथमिक भूमिका:} & \\[3pt]
Q61 & Is there one person who does most portal work for others? & {\Large$\square$} Yes, one person {\Large$\square$} A few share {\Large$\square$} Everyone does own {\Large$\square$} N/A \\
 & {\hindifont\small\itshape\color{black!70} क्या कोई एक व्यक्ति दूसरों का पोर्टल कार्य करता है?} & \\[3pt]
Q62 & Is the training appropriate for your specific job role? & \likert \\
 & {\hindifont\small\itshape\color{black!70} क्या प्रशिक्षण आपकी विशिष्ट भूमिका के लिए उपयुक्त है?} & \\[3pt]
Q63 & When a portal gives an error, what do you do? & {\Large$\square$} Wait for IT {\Large$\square$} Ask colleague {\Large$\square$} Use paper {\Large$\square$} Fix myself {\Large$\square$} Tell supervisor {\Large$\square$} Abandon \\
 & {\hindifont\small\itshape\color{black!70} जब पोर्टल त्रुटि देता है, तो आप क्या करते हैं?} & \\[3pt]
Q64 & Do senior officers use digital tools themselves? & {\Large$\square$} Yes/{\hindifont हाँ} {\Large$\square$} Rely on subordinates {\Large$\square$} Mixed {\Large$\square$} Don't know \\
 & {\hindifont\small\itshape\color{black!70} क्या वरिष्ठ अधिकारी स्वयं डिजिटल उपकरण उपयोग करते हैं?} & \\[3pt]
Q65 & Have digital tools changed the work expected at your level? & \likert{} {\scriptsize(1=No change, 5=Completely)} \\
 & {\hindifont\small\itshape\color{black!70} क्या डिजिटल उपकरणों ने आपके स्तर पर अपेक्षित कार्य बदला है?} & \\[3pt]
Q66 & Are there digital tasks beyond your current skill level? & {\Large$\square$} Yes/{\hindifont हाँ} {\Large$\square$} No/{\hindifont नहीं}\quad If yes: \rule{4cm}{0.4pt} \\
 & {\hindifont\small\itshape\color{black!70} क्या कोई डिजिटल कार्य आपके कौशल स्तर से परे हैं?} & \\[3pt]
Q67 & Should training differ based on job level? & {\Large$\square$} Yes {\Large$\square$} Somewhat {\Large$\square$} No\quad Why: \rule{3.5cm}{0.4pt} \\
 & {\hindifont\small\itshape\color{black!70} क्या प्रशिक्षण पद स्तर के अनुसार अलग होना चाहिए?} & \\[3pt]
\end{tabularx}
\vspace{4pt}
\noindent\rule{\textwidth}{0.3pt}\vspace{3pt}
\noindent{\large\bfseries Section O: Personal \& Non-Government Tools}\\
{\hindifont\itshape\color{black!70} व्यक्तिगत एवं गैर-सरकारी उपकरण}\vspace{3pt}

\noindent\begin{tabularx}{\textwidth}{@{}L{0.8cm} X L{5.2cm}@{}}
Q68 & Do you use any non-government apps or tools to help with your work? & \ymark \\
 & {\hindifont\small\itshape\color{black!70} क्या आप अपने काम में कोई गैर-सरकारी ऐप्स या उपकरण उपयोग करते हैं?} & \\[3pt]
Q69 & Which of these do you use for work-related tasks? (tick all) & {\Large$\square$} WhatsApp {\Large$\square$} Google Docs/Drive {\Large$\square$} ChatGPT/AI {\Large$\square$} YouTube {\Large$\square$} Personal email {\Large$\square$} MS Office {\Large$\square$} Google Translate {\Large$\square$} Other: \rule{2cm}{0.4pt} \\
 & {\hindifont\small\itshape\color{black!70} इनमें से कौन-से आप कार्य के लिए उपयोग करते हैं? (सभी पर टिक करें)} & \\[3pt]
Q70 & What do you mainly use these tools for? (tick all) & {\Large$\square$} Drafting {\Large$\square$} Translating {\Large$\square$} Coordinating {\Large$\square$} Learning portals {\Large$\square$} Backup {\Large$\square$} Sharing files {\Large$\square$} Other: \rule{2cm}{0.4pt} \\
 & {\hindifont\small\itshape\color{black!70} आप इन उपकरणों का मुख्य रूप से किसलिए उपयोग करते हैं? (सभी पर टिक करें)} & \\[3pt]
Q71 & How often do you use these personal tools for official work? & {\Large$\square$} Daily {\Large$\square$} Few times/week {\Large$\square$} Occasionally {\Large$\square$} Rarely {\Large$\square$} Never \\
 & {\hindifont\small\itshape\color{black!70} आप कितनी बार इन व्यक्तिगत उपकरणों का आधिकारिक काम के लिए उपयोग करते हैं?} & \\[3pt]
Q72 & Do you feel these tools fill a gap that government systems don't cover? & \likert{} {\scriptsize(1=Not at all, 5=Absolutely)} \\
 & {\hindifont\small\itshape\color{black!70} क्या ये उपकरण सरकारी प्रणालियों की कमी पूरी करते हैं?} & \\[3pt]
Q73 & Any concerns about using personal tools for official work? & \rule{\linewidth}{0.4pt}\\[3pt] && \rule{\linewidth}{0.4pt} \\
 & {\hindifont\small\itshape\color{black!70} व्यक्तिगत उपकरणों के आधिकारिक उपयोग में कोई चिंता?} & \\[3pt]
\end{tabularx}

\vspace{3pt}\noindent\rule{\textwidth}{1pt}
\begin{center}
{\large\bfseries --- END OF SCHEDULE / {\hindifont अनुसूची समाप्त} ---}\\[3pt]
\textbf{Total Questions / {\hindifont कुल प्रश्न}:} 39 + 5 + 8 + 6 = \textbf{58}\\[8pt]
\textbf{Interviewer Notes / {\hindifont साक्षात्कारकर्ता नोट्स}:}\vspace{4pt}

\rule{\textwidth}{0.4pt}\vspace{4pt}
\rule{\textwidth}{0.4pt}\vspace{4pt}
\rule{\textwidth}{0.4pt}\vspace{12pt}

\textbf{Signature / {\hindifont हस्ताक्षर}:} \rule{5cm}{0.4pt}\qquad\textbf{Date / {\hindifont दिनांक}:} \_\_\_/\_\_\_/2026
\end{center}
\end{document}
